
\documentclass[12pt]{article}

\usepackage{amsmath, amssymb}

\usepackage[margin=1in]{geometry}

\begin{document}

% =========================
% Problem 1
% =========================
\section*{Problem 1}
\textbf{(a)} 

We show that $f(x) = \sin(x^2)$ is not uniformly continuous on $[1,\infty)$.

Let $\epsilon = \tfrac{1}{2}$. We will show that for every $\delta > 0$, there exist 
$x,y \in [1,\infty)$ such that $|x-y| < \delta$ but $|f(x)-f(y)| \ge \epsilon$.

Let $\delta > 0$ be given. Choose $n \in \mathbb{N}$ large enough so that
\[
\frac{\tfrac{\pi}{2}}{\sqrt{2\pi n + \tfrac{\pi}{2}} + \sqrt{2\pi n}} < \delta.
\]

Define
\[
x = \sqrt{2\pi n + \tfrac{\pi}{2}}, \quad y = \sqrt{2\pi n}.
\]
Then $x,y \in [1,\infty)$.

Moreover,
\begin{align*}
|x-y|
&= \left| \sqrt{2\pi n + \tfrac{\pi}{2}} - \sqrt{2\pi n} \right| \\
&= \frac{\tfrac{\pi}{2}}{\sqrt{2\pi n + \tfrac{\pi}{2}} + \sqrt{2\pi n}} \\
&< \delta.
\end{align*}

However,
\begin{align*}
|f(x) - f(y)|
&= |\sin(x^2) - \sin(y^2)| \\
&= |\sin(2\pi n + \tfrac{\pi}{2}) - \sin(2\pi n)| \\
&= |1 - 0| = 1 \ge \tfrac{1}{2}.
\end{align*}

Since this holds for every $\delta > 0$, $f$ is not uniformly continuous on $[1,\infty)$.

\bigskip

\textbf{(b)} 

Let $\epsilon > 0$ and let $x,y \in [1,\infty)$. Using the identity
\[
a^3 - b^3 = (a-b)(a^2 + ab + b^2),
\]
we obtain
\[
|a-b| = \frac{|a^3 - b^3|}{|a^2 + ab + b^2|}.
\]
Let $g(x) = x^{1/3}$. Then
\begin{align*}
|g(x) - g(y)|
&= |x^{1/3} - y^{1/3}| \\
&= \frac{|x-y|}{x^{2/3} + x^{1/3}y^{1/3} + y^{2/3}}.
\end{align*}

Since $x,y \ge 1$, we have
\[
x^{2/3} \ge 1, \quad x^{1/3}y^{1/3} \ge 1, \quad y^{2/3} \ge 1,
\]
so
\[
x^{2/3} + x^{1/3}y^{1/3} + y^{2/3} \ge 3.
\]
Thus,
\[
|x^{1/3} - y^{1/3}| \le \frac{|x-y|}{3}.
\]

Choose $\delta = 3\epsilon$. Then if $|x-y| < \delta$,
\[
|g(x) - g(y)| \le \frac{|x-y|}{3} < \frac{3\epsilon}{3} = \epsilon.
\]

Therefore, $g$ is uniformly continuous on $[1,\infty)$.

% =========================
% Problem 2
% =========================
\section*{Problem 2}
\textbf{(a)} 

Consider the power series
\[
\sum_{n=1}^{\infty} \frac{(-1)^n (x-2)^n}{n \cdot 3^n}.
\]
Using the Ratio Test, compute
\begin{align*}
\lim_{n \to \infty} \left| \frac{(-1)^{n+1}(x-2)^{n+1}}{(n+1)3^{n+1}} \cdot \frac{n3^n}{(-1)^n (x-2)^n} \right|
&= \lim_{n \to \infty} \left| \frac{(-1)(x-2)n}{(n+1)3} \right| \\
&= \frac{|x-2|}{3} \cdot \lim_{n \to \infty} \frac{n}{n+1} \\
&= \frac{|x-2|}{3}.
\end{align*}

By the Ratio Test, the series converges when
\[
\frac{|x-2|}{3} < 1 \quad \Longleftrightarrow \quad |x-2| < 3.
\]
Thus, the radius of convergence is
\[
R = 3.
\]

\bigskip

\textbf{(b)} 
The center is $a = 2$, so
\[
|x-2| < 3 \quad \Longleftrightarrow \quad -1 < x < 5.
\]

\medskip

\textbf{At $x = -1$:}
\begin{align*}
\sum_{n=1}^{\infty} \frac{(-1)^n (-3)^n}{n \cdot 3^n}
&= \sum_{n=1}^{\infty} \frac{(-1)^n (-1)^n 3^n}{n \cdot 3^n} \\
&= \sum_{n=1}^{\infty} \frac{(-1)^{2n}}{n} \\
&= \sum_{n=1}^{\infty} \frac{1}{n}.
\end{align*}
This is a $p$-series with $p = 1$. Since $p \le 1$, the series diverges by the $p$-series test. Hence, the series diverges at $x=-1$.

\medskip

\textbf{At $x = 5$:}
\begin{align*}
\sum_{n=1}^{\infty} \frac{(-1)^n (3)^n}{n \cdot 3^n}
&= \sum_{n=1}^{\infty} \frac{(-1)^n}{n}.
\end{align*}
This is an alternating series. Let $a_n = \frac{1}{n}$. Then:
\begin{itemize}
\item $a_n \ge 0$ for all $n$,
\item $a_n$ is decreasing since $\frac{1}{n} \ge \frac{1}{n+1}$,
\item $\lim_{n \to \infty} a_n = 0$.
\end{itemize}
Therefore, by the Alternating Series Test, the series converges at $x=5$.

\medskip

Hence, the interval of convergence is
\[
(-1, 5].
\]

% =========================
% Problem 3
% =========================
\section*{Problem 3}
\textbf{(a)} 

Let $\epsilon > 0$. For all $x \in [0,\infty)$,
\[
|f_n(x) - f(x)| 
= \left| \left(2x + \frac{1}{\sqrt{n}}\right) - 2x \right|
= \frac{1}{\sqrt{n}}.
\]

Choose $N \in \mathbb{N}$ such that
\[
\frac{1}{\sqrt{N}} < \epsilon,
\quad \text{for example } 
N = \left\lceil \frac{1}{\epsilon^2} \right\rceil.
\]

Then for all $n \ge N$ and all $x \in [0,\infty)$,
\[
|f_n(x) - f(x)| = \frac{1}{\sqrt{n}} \le \frac{1}{\sqrt{N}} < \epsilon.
\]

Therefore, $f_n \to f$ uniformly on $[0,\infty)$.

\bigskip

\textbf{(b)} 

Let $g_n(x) = \left(2x + \frac{1}{\sqrt{n}}\right)^2$ and $g(x) = 4x^2$.

Let $x \in [0,\infty)$ be fixed and let $\epsilon > 0$. We show that there exists $N \in \mathbb{N}$ such that for all $n \ge N$,
\[
|g_n(x) - g(x)| < \epsilon.
\]

Compute:
\begin{align*}
|g_n(x) - g(x)|
&= \left| \left(2x + \frac{1}{\sqrt{n}}\right)^2 - 4x^2 \right| \\
&= \frac{4x}{\sqrt{n}} + \frac{1}{n}.
\end{align*}

Choose
\[
N_1 = \left\lceil \left(\frac{8x}{\epsilon}\right)^2 \right\rceil,
\quad
N_2 = \left\lceil \frac{2}{\epsilon} \right\rceil,
\]
and let
\[
N = \max\{N_1, N_2\}.
\]

Then for all $n \ge N$, we have
\[
\frac{4x}{\sqrt{n}} \le \frac{4x}{\sqrt{N_1}} \le \frac{\epsilon}{2},
\quad
\frac{1}{n} \le \frac{1}{N_2} \le \frac{\epsilon}{2}.
\]

Thus,
\[
|g_n(x) - g(x)|
= \frac{4x}{\sqrt{n}} + \frac{1}{n}
\le \frac{\epsilon}{2} + \frac{\epsilon}{2}
= \epsilon.
\]

Therefore, $g_n(x) \to g(x)$ for each fixed $x \in [0,\infty)$, so $g_n \to g$ pointwise.

\bigskip

\textbf{(c)} 

We show that $g_n$ does not converge uniformly to $g$ on $[0,\infty)$.

Let $\epsilon = 1$. For any $n \in \mathbb{N}$, choose
\[
x_n = \frac{\sqrt{n}}{4}.
\]
Then
\begin{align*}
|g_n(x_n) - g(x_n)|
&= \frac{4x_n}{\sqrt{n}} + \frac{1}{n} \\
&= \frac{4(\sqrt{n}/4)}{\sqrt{n}} + \frac{1}{n} \\
&= 1 + \frac{1}{n} > 1 = \epsilon.
\end{align*}

Thus, for every $n$, there exists $x_n \in [0,\infty)$ such that
\[
|g_n(x_n) - g(x_n)| \ge \epsilon.
\]

Therefore, $g_n$ does not converge uniformly to $g$ on $[0,\infty)$.

% =========================
% Problem 4
% =========================
\section*{Problem 4}
\textbf{(a)}

\textbf{Case 1:} $x = 0$. Then
\[
f_n(0) = n \cdot 0 \cdot e^{-n^2 \cdot 0^2} = 0 \quad \text{for all } n,
\]
so clearly
\[
\lim_{n \to \infty} f_n(0) = 0.
\]

\textbf{Case 2:} $x \neq 0$. Write
\[
f_n(x) = n x e^{-(n x)^2}.
\]
Let $y = n x$. Then as $n \to \infty$, $y \to \infty$, and
\[
f_n(x) = y e^{-y^2}.
\]
Now
\[
\lim_{y \to \infty} y e^{-y^2} = \lim_{y \to \infty} \frac{y}{e^{y^2}} = 0
\]
by standard exponential growth.  

Therefore, the pointwise limit is
\[
f(x) = \lim_{n \to \infty} f_n(x) = 0 \quad \text{for all } x \in \mathbb{R}.
\]

\bigskip

\textbf{(b)}

Compute the derivative of $f_n$:
\[
f_n'(x) = \frac{d}{dx} \big( n x e^{-n^2 x^2} \big)
= n e^{-n^2 x^2} + n x \cdot (-2 n^2 x) e^{-n^2 x^2} 
= n e^{-n^2 x^2} (1 - 2 n^2 x^2).
\]

Set $f_n'(x) = 0$ to find critical points:
\[
1 - 2 n^2 x^2 = 0 \quad \Longrightarrow \quad x = \frac{1}{\sqrt{2} n}.
\]

At this $x$, $f_n$ attains its maximum. Evaluate:
\[
f_n\Big(\frac{1}{\sqrt{2} n}\Big) = n \cdot \frac{1}{\sqrt{2} n} \cdot e^{-n^2 \cdot (1/2 n^2)} 
= \frac{1}{\sqrt{2}} e^{-1/2} \approx 0.43.
\]

Since
\[
\sup_{x \in [0, \epsilon]} |f_n(x) - f(x)| \ge f_n\Big(\frac{1}{\sqrt{2} n}\Big) \not\to 0,
\]
it follows that $f_n$ does not converge uniformly to $f$ on any interval containing $0$.

\bigskip

\textbf{(c)}

For $x \ge a > 0$, observe
\[
0 \le f_n(x) = n x e^{-n^2 x^2} \le n x e^{-n^2 a^2} \le n e^{-n^2 a^2} \cdot x_{\max} \le n e^{-n^2 a^2} \cdot x,
\]
but the key point is that $e^{-n^2 a^2}$ decays faster than any polynomial in $n$. In particular,
\[
\lim_{n \to \infty} n e^{-n^2 a^2} = 0.
\]

Therefore, for all $x \in [a,\infty)$,
\[
|f_n(x) - f(x)| = |f_n(x)| \le n x e^{-n^2 a^2} \to 0 \quad \text{as } n \to \infty.
\]

Since the bound is independent of $x$, it follows that
\[
\sup_{x \in [a,\infty)} |f_n(x) - f(x)| \to 0,
\]
and hence $f_n \to f$ uniformly on $[a,\infty)$ for any $a > 0$.

% =========================
% Problem 5
% =========================
\section*{Problem 5}
\textbf{(a)} 

Consider $f_n(x) = \frac{x^n}{n}$ on $[0,1]$.

For $x \in [0,1)$, we have $x^n \to 0$ as $n \to \infty$, hence
\[
f_n(x) = \frac{x^n}{n} \to 0.
\]
For $x = 1$,
\[
f_n(1) = \frac{1}{n} \to 0.
\]
Thus, the pointwise limit is
\[
f(x) = \lim_{n \to \infty} f_n(x) = 0 \quad \text{for all } x \in [0,1].
\]

\bigskip

\textbf{(b)} 

We compute
\[
\sup_{x \in [0,1]} |f_n(x) - f(x)| 
= \sup_{x \in [0,1]} \frac{x^n}{n}.
\]
Since $0 \le x \le 1$, the maximum of $x^n$ occurs at $x=1$, so
\[
\sup_{x \in [0,1]} \frac{x^n}{n} = \frac{1}{n}.
\]
Since $\frac{1}{n} \to 0$ as $n \to \infty$, it follows that
\[
\sup_{x \in [0,1]} |f_n(x) - f(x)| \to 0.
\]
Therefore, $f_n \to f$ uniformly on $[0,1]$.

\bigskip

\textbf{(c)} 

Let $h_n(x) = f_n'(x)$. Then
\[
h_n(x) = \frac{d}{dx}\left(\frac{x^n}{n}\right) = x^{n-1}.
\]

\medskip

\textbf{(i)} 

For $x \in [0,1)$, we have $x^{n-1} \to 0$, and for $x=1$,
\[
h_n(1) = 1^{n-1} = 1.
\]
Thus, the pointwise limit is
\[
h(x) =
\begin{cases}
0, & 0 \le x < 1, \\
1, & x = 1.
\end{cases}
\]

\medskip

\textbf{(ii)} 

Each $h_n$ is continuous on $[0,1]$ because it is differentiable. However, the limit function $h$ is not continuous at $x=1$.

Since the uniform limit of continuous functions must be continuous, it follows that $h_n$ does not converge uniformly to $h$ on $[0,1]$.

\end{document}
